\section{The $\kappa$ correction and the three mass versions}
\label{sec:kappa}

\paragraph{Table.} \fpath{MomentumScaleMap/run79509to79516_20261102/kappa_lookup.csv}, frozen on the cluster as
\fpath{run3pp/macro/kappa/kappa_lookup_run79509to79516_20261102.csv} (md5 \fcmd{38e0a20c...}). Columns
\fcmd{q, pT, curv, eta, phi, kappa, kappa_curv, eps, delta}; \TabRows{} nodes = \TabNq{} charges $\times$ \TabNpt{} \pt{} values
(\TabPtMin--\TabPtMax\gevc, 0.05\gevc{} steps up to 5\gevc, coarser above) $\times$ \TabNeta{} $\eta$ values (\TabEtaMin{} to
\TabEtaMax, step 0.25) $\times$ \TabNphi{} $\varphi$ values ($-\pi$ to $\pi$, step $\pi/12$; both ends present and equal).
Convention (file header): $(1/\pt)^{\rm corr}=(1/\pt)\,\kappa_{\rm curv}$ with $\kappa_{\rm curv}=1+\epsilon+q\,\delta$, and
$\kappa=1/\kappa_{\rm curv}$ is the multiplier of the transverse momentum, $\pt^{\rm corr}=\kappa\,\pt$. According to the file header the
correction is parametrised in $x=\ln\pt$ (reference \TabXref, scale \TabXscale, clamp \TabXclamp; $\kappa$ constant in $x$ outside
the clamp).
$\epsilon$ is charge-even (a scale), $q\delta$ charge-odd (a curvature offset, i.e.\ a sagitta-like distortion).

\begin{table}[H]
\centering\small
\caption{Table values $\kappa$ at $\eta=0$, $\varphi=0$ (nodes, no interpolation). Over the whole table $\kappa$ ranges from
\TabKMin{} to \TabKMax{} (extremes at large \pt{} and $|\eta|$).}
\label{tab:knodes}
\begin{tabular}{lcc}
\toprule
\pt{} (GeV/$c$) & $\kappa(q=-1)$ & $\kappa(q=+1)$ \\
\midrule
% generated by scripts/make_numbers.py
0.1 & 1.0006 & 1.0103 \\
0.2 & 0.9984 & 1.0074 \\
0.3 & 0.9952 & 1.0115 \\
0.5 & 0.9890 & 1.0211 \\
0.7 & 0.9827 & 1.0289 \\
1 & 0.9725 & 1.0378 \\
1.5 & 0.9549 & 1.0481 \\
2 & 0.9379 & 1.0555 \\
3 & 0.9074 & 1.0666 \\
5 & 0.8597 & 1.0838 \\

\bottomrule
\end{tabular}
\end{table}

The charge-odd part dominates above $\approx0.5$\gevc: at $\eta=\varphi=0$, $\kappa(q=-1)$ falls and $\kappa(q=+1)$ rises with
\pt{}. The table also varies with $\varphi$ (at $\pt=0.5$\gevc, $\eta=0$: \TabPhiLoMinus--\TabPhiHiMinus{} for $q=-1$,
\TabPhiLoPlus--\TabPhiHiPlus{} for $q=+1$) and with $\eta$ (\TabEtaLoMinus--\TabEtaHiMinus{} and \TabEtaLoPlus--\TabEtaHiPlus{}
at $\varphi=0$).

\paragraph{Application.} For each daughter of each candidate, with the stored momentum (at the fitted secondary vertex) and charge:
\begin{enumerate}
\item $\kappa$ is looked up by linear interpolation in $(\ln\pt,\eta,\varphi)$ between the eight surrounding nodes, separately
  for $q=\pm1$; outside the node range the value is clamped to the edge ($\pt<0.1$ or $>20$\gevc, $|\eta|>1.5$);
\item all three momentum components are multiplied by $\kappa$ (\pt{} scaled, $\eta$ and $\varphi$ unchanged);
\item the V0 mass is recomputed from the two daughter four-vectors with PDG masses ($m_\pi=139.570$, $m_p=938.272$\mevcc),
  and the V0 \pt{} from the sum of the daughter momenta.
\end{enumerate}

\paragraph{Mass versions.} From the same candidates and the same selection, each job writes
\begin{itemize}
\item \textbf{KFP}: the KFParticle mass and \pt{} (identical to the earlier production 480608, Sec.~\ref{sec:audit});
\item \textbf{no $\kappa$} (``raw''): mass and \pt{} recomputed from the stored daughter momenta, $\kappa\equiv1$;
\item \textbf{with $\kappa$}: as ``raw'' after the correction.
\end{itemize}

\choice{The effect of the correction is measured as ``with $\kappa$'' vs ``no $\kappa$'': same candidates, same mass formula, only
the daughter momenta differ. On run 79509 segments 0--30, recomputed minus KFP mass is \DmKsMean{} (RMS \DmKsRms)\mevcc{} for \Ks{} and
\DmLaMean{} (RMS \DmLaRms)\mevcc{} for \La{}, all candidates; the origin of the difference (KFParticle vertex fit vs stored
momenta) is not investigated.}

\paragraph{Mean correction applied.} QA-selected candidates, averaged over all masses in the histogram range (job logs):
\begin{itemize}
\item period A: \Ks{} $\pi^-$ \PerAKapKsOne, $\pi^+$ \PerAKapKsTwo; \La{} pion \PerAKapLaOne, proton \PerAKapLaTwo;
\item period B: \Ks{} $\pi^-$ \PerBKapKsOne, $\pi^+$ \PerBKapKsTwo; \La{} pion \PerBKapLaOne, proton \PerBKapLaTwo.
\end{itemize}
The \pt{} dependence on run 79509 segments 0--30 is in Table~\ref{tab:kprof}.

\begin{table}[H]
\centering\small
\caption{Mean $\kappa$ applied per daughter in 0.1\gevc{} bins of daughter \pt{} (bin containing the quoted value), all
candidates of run 79509 segments 0--30 before the QA rows (local study, \fpath{macros/fillV0MassCorr.C}); ``--'' fewer than 50
daughters. For \La{} the pion and proton columns mix both charges (\La{} and $\bar\Lambda$).}
\label{tab:kprof}
\begin{tabular}{lccccccc}
\toprule
daughter & \KprofPts \\
\midrule
% generated by scripts/make_numbers.py
\Ks: $\pi^-$ & 1.004 & 1.002 & 0.998 & 0.989 & 0.976 & 0.950 & -- \\
\Ks: $\pi^+$ & 1.010 & 1.012 & 1.019 & 1.027 & 1.035 & 1.048 & -- \\
\La+$\bar\Lambda$: $\pi$ & 1.007 & 1.007 & 1.012 & 1.019 & 1.015 & 1.029 & -- \\
\La+$\bar\Lambda$: p, $\bar{\rm p}$ & 1.008 & 1.008 & 1.008 & 1.008 & 1.004 & 0.992 & -- \\

\bottomrule
\end{tabular}
\end{table}

\caveat{The table header states that $\kappa$ was trained on vertex-constrained daughter kinematics and deployed on raw
\fcmd{SvtxTrackMap} tracks. Here it is applied to KFParticle daughter momenta constrained to the secondary vertex. Whether this is the
same constraint as in the training is not verified.}
